\section{Part I: What Few Target Labels Fix First}
\label{sec:part1-what-few-labels-fix}

Part~I asks what $K$ labels per emotion change in the linear head and how much of that change a source corpus can supply instead.

\subsection{The Gain Map}
\label{sec:part1-map}

Adding a source helps when target labels are scarce and hurts when they are plentiful (Fig.~\ref{fig:gainmap}).
For audio, the mean gain at $K=1$ is $\NUM{map-gain-a-k1}$ points, and its 95\% interval, treating the thirty iterations as independent, lies above zero for \NUM{map-gain-pos-a-k1} of the \NUM{map-pairs-a} ordered pairs.
The mean gain falls below zero by $K=25$ ($\NUM{map-gain-a-k25}$ points), and with the full pool the interval lies below zero for \NUM{map-gain-neg-a-kall} pairs.
All eight configurations follow the same course: the mean gain at $K=1$ ranges from $\NUM{map-gain-v-k1}$ (V) to $\NUM{map-gain-avt-k1}$ points (AVT), and at $K=25$ it is negative in every configuration.
The size of the gain is largely a property of the two corpora: an additive model with one effect per target and one per source explains \NUM{map-share-a-lo}--\NUM{map-share-v-hi}\% of the variance of the gains and of the largest budget at which a source still helps, across A, A~(HuBERT), V, and AV.

\begin{figure*}[t]
\centering
\FIG{gainmap}{\textwidth}
\caption{The gain map. (a)~Mean gain of pooling over target-only training (pooled minus target-only UAR in the same iteration) over ordered pairs, against $K$. (b)~For A, V, AV, and T, share of pairs whose 95\% interval of the gain lies above zero, includes zero, or lies below zero (too narrow for IEMOCAP targets, whose iterations share test folds). All ordered pairs; the pooled classifier uses the penalty of the source-only classifier.}
\label{fig:gainmap}
\end{figure*}

\subsection{Two Costs of a Few Labels}
\label{sec:part1-costs}

We assume that (A1)~within each emotion the score difference $z$ of a pair is Gaussian with a variance shared by the two emotions, so that \eqref{eq:twocloud} holds; (A2)~the learner keeps a finite $\ell_2$ penalty on the weights and none on the intercepts (with a vanishing penalty on separable data, additive corrections of the logits lose their effect~\cite{kini2021labelimbalanced}); and (A3)~the $K$ labels per emotion are a random draw from the pool.

\emph{Threshold.}
With the direction fixed, the intercept that balances two emotions places the threshold at the midpoint of the two class means of $z$.
Estimated from $K$ labels per emotion, this midpoint has variance $\sigma^2/(2K)$, so the residual offset satisfies $\mathbb{E}[o^2] = 1/(2K)$.
Expanding \eqref{eq:twocloud} to second order in $o$ gives the expected threshold loss that remains,
\begin{equation}
\label{eq:thrcost}
\mathbb{E}\left[a(d, 0) - a(d, o)\right] \approx \frac{d}{4}\,\phi\!\left(\frac{d}{2}\right)\frac{1}{2K} \le \frac{\phi(1)}{4K},
\end{equation}
where $\phi$ is the standard normal density and the maximum over $d$ is attained at $d=2$.
A handful of labels therefore removes most of any threshold offset, however large the offset was.

\emph{Direction.}
The error of an estimated direction shrinks with the ratio of an effective dimension $q$ to $K$: for a linear rule estimated from $K$ labels per emotion, the separation behaves as
\begin{equation}
\label{eq:dircost}
d(K) = \frac{\Delta}{\sqrt{1 + q/K}},
\end{equation}
the form obtained from the dimension-proportional term of the asymptotic expansion of the misclassification rate of the linear discriminant, which we read as the cost of the direction~\cite{okamoto1963asymptotic,raudys1998intrinsic}.
Fitted to the separations of target-only classifiers along the label budget, $q$ is $\NUM{p1-qdir-a}$ for WavLM audio, $\NUM{p1-qdir-ah}$ for HuBERT audio, $\NUM{p1-qdir-av}$ for AV, and $\NUM{p1-qdir-v}$ for the face, so for audio a direction comparable to the target's best needs on the order of twenty labels per emotion.
Four predictions follow; Table~\ref{tab:part1} collects the measurements that test the first two.
(P1)~Most of the loss of a source-only classifier that a few labels recover is threshold loss.
(P2)~The share of the gain that comes from the threshold falls as $K$ grows.
(P3)~Pooling corrects the threshold in proportion to the weight that the target receives (Section~\ref{sec:part1-pooling}).
(P4)~A source substitutes for a number of target labels set by $q$ and by how well its direction serves the target (Section~\ref{sec:part1-value}).

\begin{figure}[b]
\centering
\FIG{twocloud}{\columnwidth}
\caption{Threshold and direction of a pairwise decision. (a)~Two Gaussian classes along the score difference $z$; the offset $o$ moves the decision point away from the midpoint, and the direction sets the separation $d$. (b)~Threshold and direction losses of source-only classifiers against the target's best linear rule, in points of pairwise accuracy, averaged over emotion pairs and over ordered pairs from different corpora.}
\label{fig:twocloud}
\end{figure}

\subsection{What the Labels Fix}
\label{sec:part1-fix}

The two-cloud model reproduces the pairwise accuracies of the classifiers we train: over all emotion pairs and ordered pairs of the eight configurations, the accuracy predicted from the measured $d$ and $o$ correlates with the measured accuracy at \NUM{p1-e13-fit-r-lo}--\NUM{p1-e13-fit-r-hi}, with a mean absolute error of \NUM{p1-e13-fit-mae-lo}--\NUM{p1-e13-fit-mae-hi} points.

\emph{Source-only classifiers lose mainly through the direction, except on the face.}
Over ordered pairs from different corpora, the threshold loss is \NUM{p1-e13-loss-offset-a-xsame} points and the direction loss \NUM{p1-e13-loss-dir-a-xsame} points for WavLM audio (Fig.~\ref{fig:twocloud}b), \NUM{p1-e13-loss-offset-v-xsame} and \NUM{p1-e13-loss-dir-v-xsame} points for the face, and \NUM{p1-e13-loss-offset-t-xsame} and \NUM{p1-e13-loss-dir-t-xsame} points for text.
The median over pairs of the threshold's share of the loss is \NUM{p1-e29-lossshare-lo}--\NUM{p1-e29-lossshare-hi} across configurations.

\emph{The threshold loss is mostly a shift of the corpus mean.}
Centering, which subtracts the mean of the unlabeled utterances of each corpus, raises the source-only UAR by \NUM{p1-center-srconly-a} points for WavLM audio, \NUM{p1-center-srconly-ah} for HuBERT audio, and \NUM{p1-center-srconly-v} for the face, and it raises it for \NUM{p1-center-srconly-up-a-n} of the \NUM{p1-center-srconly-pairs-a-n} audio pairs (three iterations).
Refitting only the intercepts of the source-only classifier on $K=5$ labels per emotion gains \NUM{p1-d1-gain-a-k5} points for WavLM audio, close to the \NUM{p1-merged-gain-a-k5} points that pooling gains with the same labels.
After centering, the two gains shrink to \NUM{p1-center-d1-gain-a-k5} and \NUM{p1-center-merged-gain-a-k5} points (Fig.~S2 of the supplementary material).

\emph{The labels go to the threshold first (P1, P2).}
Split into a threshold part and a direction part with \eqref{eq:twocloud}, the improvement of the pooled classifier over the source-only classifier is dominated by the threshold part for single-modality audio and face at $K \le 5$: its share is \NUM{p1-e29-critshare-kle5-single-lo}--\NUM{p1-e29-critshare-kle5-single-hi} across WavLM audio, HuBERT audio, and the face at $K=1$ and $K=5$ (Fig.~\ref{fig:thresholdshare}a and Table~\ref{tab:part1}; the shares exclude the two pairs of IEMOCAP sessions).
This share is much larger than the threshold's share of the loss, so the labels correct the cheap part of the error first, and it falls with $K$ (\NUM{p1-e29-critshare-k10-ah} for HuBERT audio at $K=10$).
In the fused and text configurations the threshold loss is smaller than in audio, and at $K=5$ the threshold part is about half of the gain or less (\NUM{p1-e29-critshare-k5-fused-text-lo}--\NUM{p1-e29-critshare-k5-fused-text-hi}, with \NUM{p1-e29-critshare-k5-av} for AV); at $K=1$ it still exceeds half for AV (\NUM{p1-e29-critshare-k1-av}), VT (\NUM{p1-e29-critshare-k1-vt}), and AVT (\NUM{p1-e29-critshare-k1-avt}).
A two-layer perceptron gives threshold shares of \NUM{p1-e30-critshare-k5-audio-lo}--\NUM{p1-e30-critshare-k5-audio-hi} for audio and \NUM{p1-e30-critshare-k5-v-wd2}--\NUM{p1-e30-critshare-k5-v-wd4} for the face at $K=5$, so the pattern is not specific to the linear head, although for HuBERT audio it falls below half.

\emph{The direction stays where the source left it.}
Under \eqref{eq:dircost}, even the optimal combination of the source direction with ten labels per emotion recovers \NUM{p1-dirloss-recov-bound-k10-audio-lo}--\NUM{p1-dirloss-recov-bound-k10-audio-hi}\% of the direction loss for audio, and hypothesis transfer, which starts from the source classifier and refits it on the labels, recovers \NUM{p1-dirloss-recov-ht-k10-a}\% for WavLM and \NUM{p1-dirloss-recov-ht-k10-ah}\% for HuBERT audio (medians over pairs from different corpora).

\begin{figure}[t]
\centering
\FIG{thresholdshare}{\columnwidth}
\caption{Where the labels go. (a)~Threshold's share of the improvement of the pooled over the source-only classifier (median over ordered pairs without the two pairs of IEMOCAP sessions; ten iterations; dashed: fused and text). (b)~Expected residual errors after estimation from $K$ labels per emotion: $1/(2K)$ for the threshold offset and $q/K$ for the direction, with $q$ the median over targets.}
\label{fig:thresholdshare}
\end{figure}

\begin{table}[t]
\caption{Part~I quantities over ordered pairs from different corpora (for A~(HuBERT) the losses include all pairs): losses of the source-only classifier (points of pairwise accuracy), $q$ of \eqref{eq:dircost}, and the threshold's share of the pooled improvement (median over pairs).}
\label{tab:part1}
\centering
\footnotesize
\setlength{\tabcolsep}{3.5pt}
\begin{tabular}{@{}lrrrrrr@{}}
\toprule
 & \multicolumn{2}{c}{Loss (points)} & & \multicolumn{3}{c}{Threshold share} \\
\cmidrule(lr){2-3}\cmidrule(l){5-7}
Config. & Thr. & Dir. & $q$ & $K{=}1$ & $K{=}5$ & $K{=}10$ \\
\midrule
A & \NUM{p1-e13-loss-offset-a-xsame} & \NUM{p1-e13-loss-dir-a-xsame} & \NUM{p1-qdir-a} & \NUM{p1-e29-critshare-k1-a} & \NUM{p1-e29-critshare-k5-a} & \NUM{p1-e29-critshare-k10-a} \\
A (HuBERT) & \NUM{p1-e13-loss-offset-ah} & \NUM{p1-e13-loss-dir-ah} & \NUM{p1-qdir-ah} & \NUM{p1-e29-critshare-k1-ah} & \NUM{p1-e29-critshare-k5-ah} & \NUM{p1-e29-critshare-k10-ah} \\
V & \NUM{p1-e13-loss-offset-v-xsame} & \NUM{p1-e13-loss-dir-v-xsame} & \NUM{p1-qdir-v} & \NUM{p1-e29-critshare-k1-v} & \NUM{p1-e29-critshare-k5-v} & \NUM{p1-e29-critshare-k10-v} \\
AV & \NUM{p1-e13-loss-offset-av-xsame} & \NUM{p1-e13-loss-dir-av-xsame} & \NUM{p1-qdir-av} & \NUM{p1-e29-critshare-k1-av} & \NUM{p1-e29-critshare-k5-av} & \NUM{p1-e29-critshare-k10-av} \\
T & \NUM{p1-e13-loss-offset-t-xsame} & \NUM{p1-e13-loss-dir-t-xsame} & -- & \NUM{p1-e29-critshare-k1-t} & \NUM{p1-e29-critshare-k5-t} & \NUM{p1-e29-critshare-k10-t} \\
VT & \NUM{p1-e13-loss-offset-vt-xsame} & \NUM{p1-e13-loss-dir-vt-xsame} & -- & \NUM{p1-e29-critshare-k1-vt} & \NUM{p1-e29-critshare-k5-vt} & \NUM{p1-e29-critshare-k10-vt} \\
AT & \NUM{p1-e13-loss-offset-at-xsame} & \NUM{p1-e13-loss-dir-at-xsame} & -- & \NUM{p1-e29-critshare-k1-at} & \NUM{p1-e29-critshare-k5-at} & \NUM{p1-e29-critshare-k10-at} \\
AVT & \NUM{p1-e13-loss-offset-avt-xsame} & \NUM{p1-e13-loss-dir-avt-xsame} & -- & \NUM{p1-e29-critshare-k1-avt} & \NUM{p1-e29-critshare-k5-avt} & \NUM{p1-e29-critshare-k10-avt} \\
\bottomrule
\end{tabular}
\end{table}

\subsection{How Pooling Corrects the Threshold}
\label{sec:part1-pooling}

Pooling corrects the threshold without a separate intercept for the target; one Newton step from the source-only solution shows how (P3).
Let $N_s$ and $N_t$ be the numbers of source and target training utterances, $\kappa_s$ and $\kappa_t$ the mean curvatures of the loss at the source-only solution, $A_s$ the curvature-weighted scatter of the source features, and $\boldsymbol{\delta}$ the target mean minus the curvature-weighted source mean.
One step leaves the fraction $1 - \lambda$ of the source-only offset $o_s$, with
\begin{equation}
\label{eq:lambda}
\lambda = \frac{\rho}{1 + \rho},\qquad \rho = N_t\kappa_t\left[\frac{1}{N_s\kappa_s} + \boldsymbol{\delta}^{\top}\left(A_s + \frac{I}{C}\right)^{-1}\boldsymbol{\delta}\right].
\end{equation}
The first term is an intercept path, through which the target corrects the threshold in proportion to its share of the training weight; the second is a weight path: when the target is separable from the source along directions in which the source varies little, the pooled classifier puts weight on such a direction, which acts as an intercept that only target utterances see.
This form is the group influence of a set of points~\cite{koh2019accuracy,pregibon1981logistic} written for added points, and its intercept path reduces to the threshold shift $m/(n+m)$ of a one-dimensional discriminant pooled from $n$ target and $m$ source points~\cite{desilva2023value}; the derivation is in the supplementary material.

The measured correction $\lambda = 1 - o_m/o_s$ of pooled classifiers is ordered as \eqref{eq:lambda} predicts, although the expression overestimates its average (Fig.~S1 of the supplementary material): over ordered pairs and $K \in \{1, 5, 10\}$, the rank correlation between the predicted and the measured $\lambda$ is \NUM{p1-e15-lam-rho-lo}--\NUM{p1-e15-lam-rho-hi} across A, A~(HuBERT), V, and AV (ten iterations).
The full Newton step orders them at \NUM{p1-e16-lam-rho-lo}--\NUM{p1-e16-lam-rho-hi} and predicts the change of the direction better for V and AV (\NUM{p1-e16-dir-rho-video-lo}--\NUM{p1-e16-dir-rho-video-hi}) than for audio (\NUM{p1-e16-dir-rho-audio-lo}--\NUM{p1-e16-dir-rho-audio-hi}).
The weight path is consistent with the observation that the distance between corpora does not predict the gain: a large distance lowers the source-only accuracy but widens the path through which pooling corrects the threshold.

\subsection{What a Source Is Worth}
\label{sec:part1-value}

A source contributes a direction, whose value can be expressed in target labels (P4).
Let $e$ be the efficiency of the source direction for the target, the factor by which it scales the target's best separation, $d = e\Delta$.
Under \eqref{eq:dircost}, the direction estimated from $K$ labels has a noise-to-signal ratio $q/K$, and the source direction has the ratio $(1 - e^2)/e^2$.
Combining two independent estimates with optimal weights adds their sample sizes, so starting from the source and refitting on $K$ labels behaves like training on the target alone with $K + n$ labels per emotion, where
\begin{equation}
\label{eq:value}
n = q\,\frac{e^2}{1 - e^2}.
\end{equation}
Equation~\eqref{eq:value} inverts the classical relation between the cosine of an estimated direction and the sample size~\cite{raudys1998intrinsic,dobriban2018highdimensional}, and $n$ does not depend on $K$.

We measure $n$ by locating the UAR of hypothesis transfer on the target-only learning curve and subtracting $K$ (Fig.~S3 of the supplementary material).
For a single audio source, $n$ rises only from \NUM{p1-e32-n1src-audio-lo} to \NUM{p1-e32-n1src-audio-hi} labels per emotion while $K$ grows from $5$ to $25$; for the face it is \NUM{p1-e32-n1src-v-k5}--\NUM{p1-e32-n1src-v-k10}.
Starting from all sources pooled, $n$ is \NUM{p1-e32-nall-a-k5} and \NUM{p1-e32-nall-a-k10} labels per emotion for WavLM audio at $K=5$ and $K=10$.
Across audio pairs, the rank correlation between the predicted and the measured $n$ is \NUM{p1-e32-rho-audio-lo}--\NUM{p1-e32-rho-audio-hi}, and the measured values are \NUM{p1-e32-ratio-audio-lo}--\NUM{p1-e32-ratio-audio-hi} times the predicted ones; for all sources pooled the ratio is \NUM{p1-e32-allsrc-ratio-lo}--\NUM{p1-e32-allsrc-ratio-hi} across A, A~(HuBERT), V, and AV.

